%% =====================================================================
%% KSA / KFU - CCSIT Graduation Project Report Template (LaTeX2e)
%% THIS IS THE FILE YOU EDIT.
%% ---------------------------------------------------------------------
%% Files in this package:
%%   ccsit-preamble.tex   fixed layout            -- do not edit
%%   ccsit-macros.tex     your project details    -- EDIT THIS
%%   cover.tex            title page
%%   acknowledgment.tex   your acknowledgment     -- EDIT THIS
%%   undertaking.tex      signed declaration      -- do not edit the text
%%   bibliography.bib     your references         -- EDIT THIS
%%   CCSIT-Logo.png       college logo
%%   figures/             your images
%%
%% Compile with pdfLaTeX; Biber runs between the LaTeX passes:
%%     pdflatex main.tex / biber main / pdflatex main.tex / pdflatex main.tex
%% =====================================================================

\documentclass[11pt,a4paper]{report}

%% =====================================================================
%% KSA / KFU - CCSIT Graduation Project Report Template
%% PREAMBLE  --  DO NOT EDIT THIS FILE
%% ---------------------------------------------------------------------
%% This file holds the fixed layout of the report: packages, fonts,
%% margins, header and footer, code-listing style, and section
%% numbering. It is loaded from main.tex with %% =====================================================================
%% KSA / KFU - CCSIT Graduation Project Report Template
%% PREAMBLE  --  DO NOT EDIT THIS FILE
%% ---------------------------------------------------------------------
%% This file holds the fixed layout of the report: packages, fonts,
%% margins, header and footer, code-listing style, and section
%% numbering. It is loaded from main.tex with %% =====================================================================
%% KSA / KFU - CCSIT Graduation Project Report Template
%% PREAMBLE  --  DO NOT EDIT THIS FILE
%% ---------------------------------------------------------------------
%% This file holds the fixed layout of the report: packages, fonts,
%% margins, header and footer, code-listing style, and section
%% numbering. It is loaded from main.tex with \input{ccsit-preamble}.
%%
%% Everything you need to change is in ccsit-macros.tex, in the block headed
%% "CCSIT / KFU MACROS (EDIT THESE)".
%% =====================================================================

%% ----------------------------- Language & Fonts -----------------------------
\usepackage[english]{babel}
\usepackage{csquotes}
\usepackage[utf8]{inputenc}
\usepackage[T1]{fontenc}
\usepackage{newtxtext,newtxmath}

%% ----------------------------- Layout & Spacing -----------------------------
\usepackage[a4paper,left=25mm,right=25mm,top=25mm,bottom=25mm,headheight=14pt,footskip=12mm]{geometry}
\usepackage{setspace}
\onehalfspacing
\usepackage{indentfirst}      % indent first paragraph after headings
\setlength{\parindent}{1.5em} % indent size
\setlength{\parskip}{0pt}     % keep paragraphs tight

%% ----------------------------- Figures & Tables -----------------------------
\usepackage{graphicx}
\usepackage{float}
\usepackage{caption}
\usepackage{subcaption}
\usepackage{longtable}
\usepackage{booktabs}
\usepackage{multirow}
\usepackage[table]{xcolor}

%% ----------------------------- Math, Units, URLs -----------------------------
\usepackage{amsmath}
\usepackage{xurl}
\usepackage{siunitx}

%% ----------------------------- Code Listings -----------------------------
\usepackage{listings}
\lstset{
  basicstyle=\small\ttfamily,
  numbers=left,
  numberstyle=\scriptsize\ttfamily,
  numbersep=8pt,
  frame=tb,
  float=htb,
  aboveskip=6mm,
  belowskip=3mm,
  backgroundcolor=\color{black!2},
  showstringspaces=false,
  breaklines=true,
  tabsize=2,
  captionpos=t
}

%% ----------------------------- Header / Footer -----------------------------
%% \ccsitCollege and \gpTitleShort are defined in ccsit-macros.tex, which
%% main.tex loads after this file. That is fine: fancyhdr expands them at
%% page shipout.

\usepackage{fancyhdr}
\pagestyle{fancy}
\fancyhf{}
\lhead{\small \ccsitCollege}
\rhead{\small \textit{\gpTitleShort}}
\cfoot{\small \thepage}
\renewcommand{\headrulewidth}{0.2pt}
\renewcommand{\footrulewidth}{0pt}

%% ----------------------------- Hyperref -----------------------------
\usepackage{hyperref}
\hypersetup{
  colorlinks=true,
  linkcolor=blue!40!black,
  urlcolor=blue!40!black,
  citecolor=blue!40!black
}

%% ----------------------------- Acronyms -----------------------------
\usepackage{acronym}

%% ----------------------------- Bibliography (IEEE style) -----------------------------
\usepackage[backend=biber,style=ieee]{biblatex}
\addbibresource{bibliography.bib}

%% ----------------------------- Graphics path -----------------------------
\graphicspath{{figures/}}

%% ----------------------------- Numbered Sections (Not Chapters) -----------------------------
\makeatletter
\renewcommand{\thesection}{\arabic{section}}
\renewcommand{\thesubsection}{\thesection.\arabic{subsection}}
\renewcommand{\thesubsubsection}{\thesubsection.\arabic{subsubsection}}
\makeatother

\setcounter{secnumdepth}{3}.
%%
%% Everything you need to change is in ccsit-macros.tex, in the block headed
%% "CCSIT / KFU MACROS (EDIT THESE)".
%% =====================================================================

%% ----------------------------- Language & Fonts -----------------------------
\usepackage[english]{babel}
\usepackage{csquotes}
\usepackage[utf8]{inputenc}
\usepackage[T1]{fontenc}
\usepackage{newtxtext,newtxmath}

%% ----------------------------- Layout & Spacing -----------------------------
\usepackage[a4paper,left=25mm,right=25mm,top=25mm,bottom=25mm,headheight=14pt,footskip=12mm]{geometry}
\usepackage{setspace}
\onehalfspacing
\usepackage{indentfirst}      % indent first paragraph after headings
\setlength{\parindent}{1.5em} % indent size
\setlength{\parskip}{0pt}     % keep paragraphs tight

%% ----------------------------- Figures & Tables -----------------------------
\usepackage{graphicx}
\usepackage{float}
\usepackage{caption}
\usepackage{subcaption}
\usepackage{longtable}
\usepackage{booktabs}
\usepackage{multirow}
\usepackage[table]{xcolor}

%% ----------------------------- Math, Units, URLs -----------------------------
\usepackage{amsmath}
\usepackage{xurl}
\usepackage{siunitx}

%% ----------------------------- Code Listings -----------------------------
\usepackage{listings}
\lstset{
  basicstyle=\small\ttfamily,
  numbers=left,
  numberstyle=\scriptsize\ttfamily,
  numbersep=8pt,
  frame=tb,
  float=htb,
  aboveskip=6mm,
  belowskip=3mm,
  backgroundcolor=\color{black!2},
  showstringspaces=false,
  breaklines=true,
  tabsize=2,
  captionpos=t
}

%% ----------------------------- Header / Footer -----------------------------
%% \ccsitCollege and \gpTitleShort are defined in ccsit-macros.tex, which
%% main.tex loads after this file. That is fine: fancyhdr expands them at
%% page shipout.

\usepackage{fancyhdr}
\pagestyle{fancy}
\fancyhf{}
\lhead{\small \ccsitCollege}
\rhead{\small \textit{\gpTitleShort}}
\cfoot{\small \thepage}
\renewcommand{\headrulewidth}{0.2pt}
\renewcommand{\footrulewidth}{0pt}

%% ----------------------------- Hyperref -----------------------------
\usepackage{hyperref}
\hypersetup{
  colorlinks=true,
  linkcolor=blue!40!black,
  urlcolor=blue!40!black,
  citecolor=blue!40!black
}

%% ----------------------------- Acronyms -----------------------------
\usepackage{acronym}

%% ----------------------------- Bibliography (IEEE style) -----------------------------
\usepackage[backend=biber,style=ieee]{biblatex}
\addbibresource{bibliography.bib}

%% ----------------------------- Graphics path -----------------------------
\graphicspath{{figures/}}

%% ----------------------------- Numbered Sections (Not Chapters) -----------------------------
\makeatletter
\renewcommand{\thesection}{\arabic{section}}
\renewcommand{\thesubsection}{\thesection.\arabic{subsection}}
\renewcommand{\thesubsubsection}{\thesubsection.\arabic{subsubsection}}
\makeatother

\setcounter{secnumdepth}{3}.
%%
%% Everything you need to change is in ccsit-macros.tex, in the block headed
%% "CCSIT / KFU MACROS (EDIT THESE)".
%% =====================================================================

%% ----------------------------- Language & Fonts -----------------------------
\usepackage[english]{babel}
\usepackage{csquotes}
\usepackage[utf8]{inputenc}
\usepackage[T1]{fontenc}
\usepackage{newtxtext,newtxmath}

%% ----------------------------- Layout & Spacing -----------------------------
\usepackage[a4paper,left=25mm,right=25mm,top=25mm,bottom=25mm,headheight=14pt,footskip=12mm]{geometry}
\usepackage{setspace}
\onehalfspacing
\usepackage{indentfirst}      % indent first paragraph after headings
\setlength{\parindent}{1.5em} % indent size
\setlength{\parskip}{0pt}     % keep paragraphs tight

%% ----------------------------- Figures & Tables -----------------------------
\usepackage{graphicx}
\usepackage{float}
\usepackage{caption}
\usepackage{subcaption}
\usepackage{longtable}
\usepackage{booktabs}
\usepackage{multirow}
\usepackage[table]{xcolor}

%% ----------------------------- Math, Units, URLs -----------------------------
\usepackage{amsmath}
\usepackage{xurl}
\usepackage{siunitx}

%% ----------------------------- Code Listings -----------------------------
\usepackage{listings}
\lstset{
  basicstyle=\small\ttfamily,
  numbers=left,
  numberstyle=\scriptsize\ttfamily,
  numbersep=8pt,
  frame=tb,
  float=htb,
  aboveskip=6mm,
  belowskip=3mm,
  backgroundcolor=\color{black!2},
  showstringspaces=false,
  breaklines=true,
  tabsize=2,
  captionpos=t
}

%% ----------------------------- Header / Footer -----------------------------
%% \ccsitCollege and \gpTitleShort are defined in ccsit-macros.tex, which
%% main.tex loads after this file. That is fine: fancyhdr expands them at
%% page shipout.

\usepackage{fancyhdr}
\pagestyle{fancy}
\fancyhf{}
\lhead{\small \ccsitCollege}
\rhead{\small \textit{\gpTitleShort}}
\cfoot{\small \thepage}
\renewcommand{\headrulewidth}{0.2pt}
\renewcommand{\footrulewidth}{0pt}

%% ----------------------------- Hyperref -----------------------------
\usepackage{hyperref}
\hypersetup{
  colorlinks=true,
  linkcolor=blue!40!black,
  urlcolor=blue!40!black,
  citecolor=blue!40!black
}

%% ----------------------------- Acronyms -----------------------------
\usepackage{acronym}

%% ----------------------------- Bibliography (IEEE style) -----------------------------
\usepackage[backend=biber,style=ieee]{biblatex}
\addbibresource{bibliography.bib}

%% ----------------------------- Graphics path -----------------------------
\graphicspath{{figures/}}

%% ----------------------------- Numbered Sections (Not Chapters) -----------------------------
\makeatletter
\renewcommand{\thesection}{\arabic{section}}
\renewcommand{\thesubsection}{\thesection.\arabic{subsection}}
\renewcommand{\thesubsubsection}{\thesubsection.\arabic{subsubsection}}
\makeatother

\setcounter{secnumdepth}{3}
%% =====================================================================
%%                 CCSIT / KFU MACROS (EDIT THESE)
%% =====================================================================
\newcommand{\ksa}{Kingdom of Saudi Arabia}
\newcommand{\moe}{Ministry of Education}
\newcommand{\kfu}{King Faisal University}
\newcommand{\ccsitCollege}{CCSIT, King Faisal University}

\newcommand{\ccsitFull}{College of Computer Sciences \& Information Technology}

\newcommand{\gpTitle}{AI Based Research Documents Repository Using Local Language Models}
\newcommand{\gpTitleShort}{[Short Title]} % used in header (keep short)
\newcommand{\gpDegree}{Bachelor of Science in Computer Science}
\newcommand{\gpDepartment}{Department of CS}

\newcommand{\gpSemester}{Second Semester}
\newcommand{\gpAcademicYear}{2025--2026}

%% Automatic Month Year based on compile date
\newcommand{\gpSubmissionMonth}{%
\ifcase\month\or
January\or February\or March\or April\or May\or June\or
July\or August\or September\or October\or November\or December%
\fi\space \number\year}

%% Students. For a team of fewer than four, do NOT delete or blank these definitions. Delete the matching lines from the "Prepared by" table in
%% cover.tex and from the signature table in main.tex instead.
\newcommand{\gpStudentA}{Raihanah Nezar Al Ali}
\newcommand{\gpStudentAID}{223007420}
\newcommand{\gpStudentB}{Maryam Ali Alkathim}
\newcommand{\gpStudentBID}{223021603}
\newcommand{\gpStudentC}{Fadak Hani Alibrahim}
\newcommand{\gpStudentCID}{223022618}

%% Supervision / Committee
\newcommand{\gpSupervisor}{Nora Abdullah Alkhaldi}
\newcommand{\gpCoSupervisor}{[Co-Supervisor Name (if any)]}
\newcommand{\gpCommitteeA}{[Committee Member 1]}
\newcommand{\gpCommitteeB}{[Committee Member 2]}
\newcommand{\gpCommitteeC}{[Committee Member 3]}
%% ----------------------------- Acknowledgment Text -----------------------------
\newcommand{\gpAcknowledgmentText}{%
\textbf{(Replace the text below with your own acknowledgment.)}\par\vspace{2mm}
We would like to express our sincere gratitude to our supervisor, \textbf{\gpSupervisor}, for the continuous
guidance, valuable feedback, and support throughout the development of this graduation project. We
also extend our appreciation to the committee members and all faculty and staff in \ccsitFull\ for
their assistance and encouragement.
}
%% ----------------------------- Undertaking / Declaration Text -----------------------------
\newcommand{\gpUndertakingText}{%
This is to declare that the Graduation Project entitled \textbf{``\gpTitle''} is an original work carried out by the undersigned, in partial fulfillment of the requirements for the degree \textbf{\gpDegree} at the \gpDepartment, \ccsitFull, \kfu.

All the analysis, design, implementation, and system development activities have been accomplished by the undersigned. Furthermore, this project has not been submitted, in whole or in part, to any other college or university for any degree or qualification.
}

%% =====================================================================
%%                           DOCUMENT
%% =====================================================================
\begin{document}

%% ----------------------------- Title Page -----------------------------
%% =====================================================================
%% cover.tex  --  the title page
%% ---------------------------------------------------------------------
%% Built entirely from the macros in ccsit-macros.tex. To change what
%% appears here, edit those macros.
%%
%% Two exceptions you may edit in this file:
%%   - teams of fewer than four students: delete the unused lines from
%%     the "Prepared by" table (and from undertaking.tex);
%%   - to show the co-supervisor or committee members 2 and 3, remove
%%     the leading % from those lines.
%% =====================================================================

\begin{titlepage}
\thispagestyle{empty}
\begin{center}

\vspace{4mm}
\includegraphics[width=0.5\textwidth]{CCSIT-Logo.png}

\vspace{6mm}

{\Large \textbf{\ksa}}\\[2mm]
{\Large \textbf{\moe}}\\[2mm]
{\Large \textbf{\kfu}}\\[2mm]
{\Large \textbf{\ccsitCollege}}\\[10mm]
{\large \textbf{\ccsitFull}}\\[10mm]
\vfill

{\LARGE \textbf{\gpTitle}\par}
\vspace{6mm}

{\large\itshape
\setstretch{1.05}
A project submitted\par
in partial fulfillment of the requirements for the degree of\par
\gpDegree\par
}

\vfill

\textbf{Prepared by}\\[2mm]
\begin{tabular}{rl}
\gpStudentA & (\gpStudentAID) \\
\gpStudentB & (\gpStudentBID) \\
\gpStudentC & (\gpStudentCID) \\
\end{tabular}

\vfill

\textbf{Supervised by}\\[2mm]
\gpSupervisor\\
% \gpCoSupervisor

\vspace{6mm}

\textbf{Committee Members}\\[2mm]
\gpCommitteeA\\
% \gpCommitteeB\\
% \gpCommitteeC\\

\vfill

{\gpSubmissionMonth}

\end{center}
\end{titlepage}

%% ----------------------------- Acknowledgment Page -----------------------------
\cleardoublepage
\thispagestyle{empty}
\begin{center}
{\Large \textbf{Acknowledgment}}
\end{center}

\vspace{6mm}
\gpAcknowledgmentText

%% ----------------------------- Undertaking -----------------------------
\cleardoublepage
\thispagestyle{empty}
\begin{center}
{\Large \textbf{Undertaking}}
\end{center}

\vspace{6mm}
\gpUndertakingText

\vfill
\noindent \textbf{Prepared by:}\\[3mm]
\begin{tabular}{@{}p{0.62\textwidth}p{0.34\textwidth}@{}}
\gpStudentA\ (\gpStudentAID) & Signature:\ \rule{45mm}{0.4pt}\\[4mm]
\gpStudentB\ (\gpStudentBID) & Signature:\ \rule{45mm}{0.4pt}\\[4mm]
\gpStudentC\ (\gpStudentCID) & Signature:\ \rule{45mm}{0.4pt}\\
\end{tabular}

\vspace{6mm}
\noindent Date:\ \rule{55mm}{0.4pt}

%% ----------------------------- Front Matter -----------------------------
\cleardoublepage
\pagenumbering{roman}
\setcounter{page}{1}

\section*{Abstract}
\addcontentsline{toc}{section}{Abstract}

% Write a concise abstract (250--300 words).
[Write your abstract here (250--300 words).]

\vspace{4mm}
\noindent\textbf{Keywords:} keyword1; keyword2; keyword3; keyword4

\cleardoublepage
\tableofcontents

\cleardoublepage
\listoftables
\addcontentsline{toc}{section}{List of Tables}




\cleardoublepage
\listoffigures
\addcontentsline{toc}{section}{List of Figures}

%% ----------------------------- Acronyms -----------------------------
\cleardoublepage
\section*{List of Acronyms}
\addcontentsline{toc}{section}{List of Acronyms}

\begin{acronym}[123456]
  \acro{API}{Application Programming Interface}
  \acro{AI}{Artificial Intelligence}
  \acro{NLP}{Natural Language Processing}
  \acro{GPU}{Graphics Processing Unit}
\end{acronym}

%% ----------------------------- Main Matter -----------------------------
\cleardoublepage
\pagenumbering{arabic}
\setcounter{page}{1}

\section{First-Heading}
\subsection{First-Sub-heading-1}
\subsubsection{First-Sub-subheading-1-a}
\subsubsection{First-Sub-subheading-1-b}
\subsection{First-Sub-heading-2}
\subsubsection{First-Sub-subheading-2-a}
\subsubsection{First-Sub-subheading-2-b}

\section{Second-Heading}
\subsection{Second-Sub-heading-1}
\subsubsection{Second-Sub-subheading-1-a}
\subsubsection{Second-Sub-subheading-1-b}
\subsection{Second-Sub-heading-2}
\subsubsection{Second-Sub-subheading-2-a}
\subsubsection{Second-Sub-subheading-2-b}
Hello World! This project follows software engineering principles \cite{Tanaka2025VDocRAGRG}.
Similar AI-based solutions have been reported in recent studies \cite{Dong2026DocResearcherAU}.
A comparable system design was presented in \cite{Gosmar2026MADPAM}.

Acronyms are expanded automatically on first use. 
The proposed system exposes an \ac{API} that applies
\ac{AI} and \ac{NLP} techniques on a \ac{GPU}, and every later mention of the \ac{API} prints the
short form only.

%% ----------------------------- References -----------------------------
\cleardoublepage
\printbibliography
\addcontentsline{toc}{section}{References}

%% ----------------------------- Appendices -----------------------------
\appendix
\cleardoublepage
%% ----------------------------- AppendixA -----------------------------
\section*{Appendix A: Additional Materials}
\addcontentsline{toc}{section}{Appendix A: Additional Materials}

The following equation is provided only as an illustrative example to demonstrate how to include a general numbered equation in this report template.

\begin{equation}
\label{eq:general-example}
y = ax + b
\end{equation}

\begin{subequations}
\label{eq:subequations-example}
\begin{align}
y &= ax + b \label{eq:subeq-a}\\
x &= \frac{y-b}{a} \label{eq:subeq-b}\\
x^{10} &= y - b \label{eq:subeq-c}\\
f^{(k)}(x) &\neq f_i(x)
\end{align}
\end{subequations}

Table~\ref{tab:general-comparison} is included only as an example to demonstrate table formatting in this report template.

\begin{table}[H]
\centering
\caption{General comparison between existing solutions and the proposed system.}
\label{tab:general-comparison}
\begin{tabular}{|l|c|c|}
\hline
\textbf{Criteria} & \textbf{Existing Solutions} & \textbf{Proposed System} \\
\hline
Ease of use & Moderate & High \\
\hline
Accuracy of results & Moderate & Improved \\
\hline
User interface clarity & Basic & Enhanced \\
\hline
Automation level & Limited & High \\
\hline
Reliability & Moderate & High \\
\hline
Scalability & Limited & Good \\
\hline
Reporting and documentation & Basic & Detailed \\
\hline
\end{tabular}
\end{table}

The following table is provided only as an illustrative example to demonstrate how to include a bordered table with \textbf{merged column headers} in this report template (Table~\ref{tab:general-comparison-merged}).

\begin{table}[H]
\centering
\caption{Comparison of system characteristics.}
\label{tab:general-comparison-merged}
\begin{tabular}{|l|c|c|}
\hline
\textbf{Feature} & \multicolumn{2}{c|}{\textbf{System Type}} \\
\cline{2-3}
 & \textbf{Traditional Approach} & \textbf{Proposed Approach} \\
\hline
User experience & Average & Improved \\
\hline
System flexibility & Low & High \\
\hline
Maintenance effort & High & Lower \\
\hline
Processing speed & Moderate & Faster \\
\hline
Overall efficiency & Moderate & High \\
\hline
\end{tabular}
\end{table}

The following table is provided only as an illustrative example to demonstrate how to include a bordered table with \textbf{merged rows} in this report template (Table~\ref{tab:general-comparison-categories}).

\begin{table}[H]
\centering
\caption{General comparison of project features.}
\label{tab:general-comparison-categories}
\begin{tabular}{|l|l|c|c|}
\hline
\textbf{Category} & \textbf{Criteria} & \textbf{Existing} & \textbf{Proposed} \\
\hline
\multirow{2}{*}{Usability}
& Ease of use & Medium & High \\
\cline{2-4}
& Accessibility & Medium & High \\
\hline
\multirow{2}{*}{Performance}
& Reliability & Medium & High \\
\cline{2-4}
& Response speed & Medium & High \\
\hline
\multirow{2}{*}{Output Quality}
& Report clarity & Basic & Detailed \\
\cline{2-4}
& Output consistency & Medium & High \\
\hline
\end{tabular}
\end{table}

The following table is provided only as an illustrative example to demonstrate how to include a bordered table with \textbf{merged rows and merged column headers} in this report template (Table~\ref{tab:merged-rows-cols-example}).

\begin{table}[H]
\centering
\caption{Illustrative comparison table (example only).}
\label{tab:merged-rows-cols-example}
\begin{tabular}{|l|l|c|c|}
\hline
\multirow{2}{*}{\textbf{Category}}
& \multirow{2}{*}{\textbf{Criteria}}
& \multicolumn{2}{c|}{\textbf{Comparison}} \\
\cline{3-4}
 &  & \textbf{Option A} & \textbf{Option B} \\
\hline
\multirow{2}{*}{Usability}
& Ease of use & High & Medium \\
\cline{2-4}
& Clarity of interface & High & Medium \\
\hline
\multirow{2}{*}{Quality}
& Reliability & High & Medium \\
\cline{2-4}
& Output clarity & High & Medium \\
\hline
\multirow{2}{*}{Efficiency}
& Time required & Low & Medium \\
\cline{2-4}
& Effort required & Low & Medium \\
\hline
\end{tabular}
\end{table}

The following code listing is provided only as an illustrative example to demonstrate how to include a numbered and captioned code excerpt in this report template (Listing~\ref{lst:example}).

\begin{lstlisting}[language=Python, caption={Example function that normalizes an input string.}, label={lst:example}]
def normalize(text):
    """Return the text in lower case, without leading or trailing spaces."""
    return text.strip().lower()
\end{lstlisting}

As shown in Listing~\ref{lst:example}, source code is typeset with line numbers and a caption placed above the excerpt. Set \texttt{language} to the language of your own code, for example \texttt{Java}, \texttt{C}, \texttt{Python}, or \texttt{SQL}. Include only short, meaningful excerpts in the report; complete source files belong in the submitted code archive.

%% ----------------------------- AppendixB -----------------------------
\section*{Appendix B: Project Screenshots}
\addcontentsline{toc}{section}{Appendix B: Project Screenshots}

\begin{figure}[H]
\centering
\includegraphics[width=0.5\textwidth]{ex.png}
\caption{Official CCSIT logo used in this graduation project report.}
\label{fig:ccsit-logo-overall}
\end{figure}

As shown in Figure~\ref{fig:ccsit-logo-overall}, the official CCSIT logo is presented for institutional identification and report branding.

\begin{figure}[H]
\centering

\begin{subfigure}[b]{0.45\textwidth}
\centering
\includegraphics[width=0.75\textwidth]{ex.png}
\caption{King Faisal University (KFU) logo.}
\label{fig:kfu-logo}
\end{subfigure}
\hfill
\begin{subfigure}[b]{0.45\textwidth}
\centering
\includegraphics[width=0.75\textwidth]{ex.png}
\caption{College of Computer Sciences \& Information Technology (CCSIT) logo.}
\label{fig:ccsit-logo}
\end{subfigure}

\caption{Institutional logos used for official report branding.}
\label{fig:logos}
\end{figure}

Figure~\ref{fig:logos} presents the institutional logos used in this report.

\end{document}